\documentclass[11pt]{article}
\usepackage[margin=1in]{geometry}
\usepackage{amsmath,amssymb,amsthm,mathtools}
\usepackage{graphicx}
\usepackage{microtype}
\usepackage[hidelinks]{hyperref}

\newtheorem{theorem}{Theorem}
\newtheorem{lemma}[theorem]{Lemma}
\newcommand{\norm}[1]{\left\lVert#1\right\rVert}
\newcommand{\R}{\mathbb R}
\newcommand{\cF}{\mathcal F}

\title{The critical $q=1$ threshold for a full derivative\\
of velocity averaging is sharp}
\author{Candidate counterexample for arXiv:1206.6659}
\date{12 August 2026}

\begin{document}
\maketitle

\noindent\textbf{Status.}
Candidate full counterexample, likely valid; human review is needed.

\section{Source question}

Ars\'enio and Masmoudi prove a full gain of $1+b-a$ derivatives at the
critical index
\[
 \beta+D(1-1/p)=1
\]
when the Besov summation exponent is $q=1$.  In Theorem~3.2 they obtain only
$s-\varepsilon$ when $q\ne1$, and immediately ask whether the full gain is
nevertheless possible \cite{ArsenioMasmoudi}.

\begin{figure}[ht]
 \centering
 \includegraphics[width=\textwidth]{figures/open_problem_crop.png}
 \caption{The open critical-summation question in
 \cite{ArsenioMasmoudi}, PDF page 8.}
\end{figure}

The following construction answers the question negatively for every
$q\in(1,\infty]$.  It uses the admissible special case
\[
 D=p=1,\qquad \beta=1,\qquad a=b=0,\qquad \alpha=\tfrac12. \tag{1.1}
\]
For these parameters, a full gain would assert
\[
 \left.
 \begin{array}{c}
 f\in B^{0,1/2}_{1,1,q}(\R_x\times\R_v),\\
 g=v\partial_xf\in B^{0,1}_{1,1,q}(\R_x\times\R_v)
 \end{array}\right\}
 \quad\Longrightarrow\quad
 \rho(x):=\int_\R f(x,v)\,dv\in B^1_{1,q}(\R_x).          \tag{1.2}
\]
The mixed spaces and Littlewood--Paley blocks below are those defined in the
source paper.

\section{Counterexample theorem}

\begin{theorem}\label{thm:counterexample}
For every $q\in(1,\infty]$ there are real tempered distributions $f,g$,
supported in a fixed compact velocity interval, such that
\[
 f\in B^{0,1/2}_{1,1,q},\qquad
 g=v\partial_xf\in B^{0,1}_{1,1,q},                       \tag{2.1}
\]
but
\[
 \int_\R f(x,v)\,dv\notin B^1_{1,q}(\R).                 \tag{2.2}
\]
The velocity average is nevertheless an $L^1$ function.  Hence (1.2) is
false, and the critical $q=1$ hypothesis in the source's full-gain theorem is
sharp within this Besov scale.
\end{theorem}

\section{Proof intuition}

The obstruction is exactly the endpoint sequence-space sum visible in the
source proof.  At spatial frequency $\lambda_m$, put $m$ shrinking velocity
atoms into $f$.  Multiplication by $v$ gives the atom at scale $2^{-j}$ an
extra factor $2^{-j}$; after one velocity derivative this makes the
$B^1_{1,q}$ transport cost of the $m$ atoms comparable to their $\ell^q$
coefficient norm.  Normalizing by $m^{-1/q}$ keeps that cost constant.

Velocity averaging behaves differently: all atoms have the same nonzero
integral and therefore add in $\ell^1$.  Their average has size
$m\,m^{-1/q}=m^{1/q'}$.  A final spatial coefficient sequence is in
$\ell^q$, so both hypotheses hold, but multiplication by $m^{1/q'}$ takes it
out of $\ell^q$.  When $q=1$, $q'=\infty$ and the bad factor disappears,
exactly matching the positive theorem.

\section{Velocity atoms}

Choose $0\le\psi\in C_c^\infty((1,3/2))$ with
\[
 c_\psi:=\int_\R\psi(v)\,dv>0,
\]
and set, for $j\ge1$,
\[
 \psi_j(v)=2^j\psi(2^jv),\qquad
 a_j(v)=v\psi_j(v)=a(2^jv),\quad a(w)=w\psi(w).           \tag{4.1}
\]
Thus
\[
 \int\psi_j=c_\psi,qquad \norm{a_j}_{L^1}=2^{-j}\norm a_{L^1}. \tag{4.2}
\]

\begin{lemma}[Critical atomic estimate]\label{lem:atoms}
For $1<q\le\infty$ and every finite scalar sequence $(c_j)$,
\[
 \norm{\sum_jc_ja_j}_{B^1_{1,q}(\R)}
 \le C_q\norm{(c_j)}_{\ell^q}.                            \tag{4.3}
\]
Moreover,
\[
 \norm{\psi_j}_{B^{1/2}_{1,q}}\le C_q2^{j/2}.             \tag{4.4}
\]
\end{lemma}

\begin{proof}
Let $\Delta_\ell$ denote a velocity block at frequency $2^\ell$.
For $\ell\le j$, $L^1$ boundedness of the block gives
\[
 2^\ell\norm{\Delta_\ell a_j}_{L^1}\le C2^{-(j-\ell)}.  \tag{4.5}
\]
For $\ell>j$, use an integer $M\ge3$, the vanishing moments of the
high-frequency kernel, and
$\norm{\partial^Ma_j}_{L^1}=C_M2^{j(M-1)}$ to obtain
\[
 2^\ell\norm{\Delta_\ell a_j}_{L^1}
 \le C_M2^{-(M-1)(\ell-j)}.                               \tag{4.6}
\]
The two right sides form an $\ell^1$ convolution kernel in $\ell-j$.
Young's sequence inequality proves the high-block part of (4.3).  Its low
block is bounded by
$C\sum_j|c_j|2^{-j}\le C_q\norm c_{\ell^q}$.
Finally (4.4) is the standard scaling estimate (or follows from the same
two-case block calculation) for $2^j\psi(2^j\cdot)$.
\end{proof}

\section{Construction and input norms}

Choose a real Schwartz function $H$ whose Fourier transform is supported in
a narrow annulus on which one fixed Littlewood--Paley multiplier equals one.
Put
\[
 k_m=4m,\qquad \lambda_m=2^{k_m},\qquad
 H_m(x)=\lambda_mH(\lambda_mx).                           \tag{5.1}
\]
The selected spatial annuli are disjoint,
$\norm{H_m}_{L^1}=\norm H_{L^1}$, and
$\norm{H_m'}_{L^1}=\lambda_m\norm{H'}_{L^1}$.

For finite $q$, use the usual conjugate exponent $q'$; for $q=\infty$ set
$1/q=0$ and $1/q'=1$.  Define
\[
 F_m(v)=\lambda_m^{-1}m^{-1/q}\sum_{j=1}^m\psi_j(v).      \tag{5.2}
\]
Let $(d_m)$ be specified below and set
\begin{align}
 f(x,v)&=\sum_{m\ge1}d_mH_m(x)F_m(v),                    \tag{5.3}\\
 g(x,v)&=\sum_{m\ge1}d_mH_m'(x)vF_m(v).                  \tag{5.4}
\end{align}
The series converge in $\mathcal S'$ (and by norm for finite $q$), and
termwise differentiation gives $v\partial_xf=g$.

By (4.4) and the triangle inequality,
\[
 \norm{F_m}_{B^{1/2}_{1,q}}
 \le C\lambda_m^{-1}m^{-1/q}\sum_{j=1}^m2^{j/2}
 \le C\lambda_m^{-1}m^{-1/q}2^{m/2}.                    \tag{5.5}
\]
Because $\lambda_m=2^{4m}$, the right side decays exponentially.  Finite
overlap of the spatial blocks therefore gives
\[
 \norm f_{B^{0,1/2}_{1,1,q}}<\infty                       \tag{5.6}
\]
for every polynomially bounded $(d_m)$.

On the other hand, Lemma~\ref{lem:atoms} gives
\[
 \norm{vF_m}_{B^1_{1,q}}
 \le C\lambda_m^{-1}m^{-1/q}
        \norm{(1,\ldots,1)}_{\ell^q}
 =C\lambda_m^{-1}.                                      \tag{5.7}
\]
Combining this with $\norm{H_m'}_1=C\lambda_m$ and the separated spatial
annuli yields
\[
 \norm g_{B^{0,1}_{1,1,q}}
 \le C\norm{(d_m)}_{\ell^q},                             \tag{5.8}
\]
with the evident supremum interpretation when $q=\infty$.

\section{The average fails at exactly one derivative}

Every atom in (5.2) has the same integral, and hence
\[
 \int_\R F_m(v)\,dv
 =c_\psi\lambda_m^{-1}m^{1/q'}.                          \tag{6.1}
\]
Therefore
\[
 \rho(x)=c_\psi\sum_{m\ge1}
 d_m\lambda_m^{-1}m^{1/q'}H_m(x).                        \tag{6.2}
\]
The multiplier selected in (5.1) is one on the $m$th Fourier support.  Thus,
for $q<\infty$,
\[
 \norm\rho_{B^1_{1,q}}^q
 \ge c\sum_{m\ge1}|d_m|^qm^{q/q'},                      \tag{6.3}
\]
and for $q=\infty$ the right side is replaced by
$c\sup_m|d_m|m$.

If $1<q<\infty$, take
\[
 r=\frac{1+1/q}{2},\qquad d_m=m^{-r}.                    \tag{6.4}
\]
Then $rq=(q+1)/2>1$, so $(d_m)\in\ell^q$ and (5.8) is finite.  But (6.3)
contains
\[
 \sum_{m\ge1}m^{q-1-rq}=\sum_{m\ge1}m^{(q-3)/2},         \tag{6.5}
\]
which diverges for every $q>1$.  If $q=\infty$, take $d_m=1$; (5.8) is
finite while $\sup_m m=\infty$.  This proves (2.1)--(2.2).

Finally, (6.2) converges absolutely in $L^1$ because
$\lambda_m^{-1}m^{1/q'}|d_m|$ is summable.  The velocity supports in (4.1)
are all contained in a fixed compact interval.  This completes the proof of
Theorem~\ref{thm:counterexample}. \qed

\section{Verification, scope, and novelty}

The parameter check is exact:
\[
 \alpha=\tfrac12>0=-D(1-1/p),\quad b=0\ge a-1=-1,
 \quad\beta+D(1-1/p)=1.
\]
The construction fails to contradict the known $q=1$ result for exactly the
right reason: then $m^{1/q'}=1$, so an $\ell^1$ input sequence cannot make the
output sum diverge.  The nonessential script included in the packet checks
finite partial sums of (5.8) and (6.5) for representative $q$.

Cheap run indexes and bounded searches through 12 August 2026 used the arXiv
id, exact title, quoted question, ``critical velocity averaging'', ``full
gain'', ``Besov $q$'', ``sharpness'', and ``counterexample''.  They found the
source/published paper and later work in other velocity-averaging regimes,
but no later paper explicitly resolving this exact critical mixed-Besov
summation question.  The proof is self-contained beyond the source's standard
Littlewood--Paley definitions.

\section{Human-review recommendation}

Verify (4.5)--(4.6) and the finite-overlap passage from the tensor products to
the mixed Besov norms in (5.6)--(5.8).  Those are the only substantive
bookkeeping points.  Once accepted, the packet gives a full negative answer
to the source question for every $q>1$, and establishes sharpness of $q=1$.

\begin{thebibliography}{9}
\bibitem{ArsenioMasmoudi}
D. Ars\'enio and N. Masmoudi,
\emph{A new approach to velocity averaging lemmas in Besov spaces},
J. Math. Pures Appl. 101 (2014), 495--551; arXiv:1206.6659.
\end{thebibliography}

\end{document}
